\documentclass[journal]{IEEEtran}

% ---------- Packages ----------
\usepackage{amsmath}
\usepackage{amsfonts}
\usepackage{cite}
\usepackage{bm}
\usepackage[dvipsnames]{xcolor}
\usepackage{graphicx}
\usepackage{tabularx}
\usepackage{pifont}

% ---------- Author editing macros ----------
\newcommand{\GO}[1]{\textcolor{magenta}{{#1}}} 
\newcommand{\TC}[1]{\textcolor{cyan}{{#1}}} 
\newcommand{\KD}[1]{\textcolor{ForestGreen}{{#1}}}
\newcommand{\rem}[1]{\textcolor{Yellow}{{#1}}}
\newcommand{\todo}[1]{\textcolor{red}{{#1}}}
\newcommand{\disc}[1]{\textcolor{orange}{#1}}

% ---------- Title & Authors ----------
\title{A Comprehensive Tutorial and Survey on Sensor Fusion Techniques and Applications}

\author{Koullis, Themis, and GabrieliniSusini%
\thanks{The authors are with the Department of Engineering, Some University, City, Country. Correspondence e-mail: author@example.com.}%
}

\begin{document}

\maketitle

\markboth{IEEE Transactions on XXX,~Vol.~X, No.~X, Month~2024}%
{Koullis \MakeLowercase{\textit{et al.}}: A Comprehensive Tutorial and Survey on Sensor Fusion Techniques and Applications}

% ---------- Abstract ----------
\begin{abstract}
Sensor fusion is a critical aspect of modern technology, offering enhanced accuracy, reliability, and efficiency across various domains such as autonomous vehicles, healthcare, and industrial automation. This paper presents a comprehensive survey of classical and advanced sensor fusion techniques, providing an in-depth tutorial on their implementation, along with a review of state-of-the-art applications and challenges. Our aim is to bridge the gap between theoretical concepts and practical applications, helping researchers and practitioners understand the evolution and future potential of sensor fusion.
\end{abstract}

\begin{IEEEkeywords}
Sensor fusion, Kalman filter, machine learning, deep learning, multi-sensor systems, SLAM, IoT.
\end{IEEEkeywords}

% ============================================
\section{Introduction}
% ============================================

\subsection{Motivation and Importance of Sensor Fusion}

Sensor fusion has become a cornerstone of modern technological advancements across various domains, including autonomous systems, healthcare, robotics, Internet of Things (IoT), and industrial automation. With the growing complexity and interconnectedness of systems, relying on a single sensor or data source often leads to limitations in performance, reliability, and accuracy. Sensor fusion addresses these limitations by integrating data from multiple sensors, offering several key benefits:
\begin{itemize}
    \item \textbf{Enhanced Accuracy}: Combining data from different types of sensors allows systems to achieve more precise estimations and measurements than using individual sensors alone.
    \item \textbf{Improved Robustness}: Sensor fusion provides redundancy by using multiple sensors to capture the same or similar information, enabling systems to remain functional even if one or more sensors fail or provide noisy data.
    \item \textbf{Increased Completeness}: By integrating different modalities of data (e.g. spatial, temporal, spectral), sensor fusion can provide a more comprehensive understanding of the environment or subject.
    \item \textbf{Context Awareness}: In dynamic and uncertain environments, such as autonomous driving, the ability to fuse information from diverse sensors, such as LiDAR, cameras, radar, and GPS, enables more reliable decision-making.
\end{itemize}

These advantages make sensor fusion indispensable for critical applications. For instance, in autonomous vehicles, sensor fusion is used to perceive the environment, localise the vehicle, and make decisions in real time. In healthcare, wearable devices integrate signals from multiple sensors (e.g. ECG, accelerometers) to provide accurate monitoring of patients' vital signs. In industrial automation, fusing sensor data from cameras, gyroscopes, and other sources ensures high-precision control over manufacturing processes.

\subsection{Goals of This Paper}

The growing interest and research in sensor fusion have led to significant advancements in both traditional and AI-based techniques. However, many researchers, engineers, and practitioners face the challenge of navigating the vast array of available techniques, understanding their applications, and learning how to implement them effectively.

This paper aims to address these challenges by offering a comprehensive survey and tutorial on sensor fusion. It combines a structured review of classical and advanced sensor fusion methods with practical implementation examples. Our overarching goal is divided into the following objectives:
\begin{itemize}
    \item \textbf{Provide the foundations of sensor fusion:} We first provide a comprehensive tutorial on the fundamental principles behind sensor fusion techniques.
    \item \textbf{Bridge the Gap Between Theory and Practice}: We provide a detailed explanation of practical sensor fusion techniques, from traditional methods like Bayesian filtering and statistical fusion to more recent approaches, including machine learning and deep learning.
    \item \textbf{Highlight State-of-the-Art Research}: We review recent research contributions to sensor fusion, exploring cutting-edge applications such as multimodal deep learning, simultaneous localisation and mapping (SLAM), and their integration in robotics, autonomous vehicles, and healthcare.
    \item \textbf{Identify Open Challenges}: Finally, we discuss ongoing challenges in sensor fusion, such as real-time processing, scalability, sensor reliability, and uncertainty management, offering insights into promising research directions for the future.
\end{itemize}

By providing this combination of theoretical foundations, practical examples, and a review of the state of the art, this paper serves as both a valuable resource for those new to sensor fusion and an up-to-date reference for researchers and professionals in the field.

\subsection{\TC{Motivation for this tutorial and survey paper}}
\KD{[DELK: write, add other tutorials and/or surveys on sensor fusion, explain why this is timely and important!]}

Sensor fusion has long been a central topic in signal processing, robotics, control, and artificial intelligence, enabling robust state estimation and decision making by combining information from multiple heterogeneous sources. Foundational work in multisensor data fusion has established theoretical frameworks and algorithmic principles for combining information across sensing modalities, with early tutorials and surveys providing comprehensive overviews of probabilistic, model based, and architectural approaches \cite{hall_introduction_1997, luo_multisensor_1989}. These works, together with later consolidated treatments such as the multisensor data fusion handbook \cite{liggins_handbook_2009}, form the basis of much of the classical literature in the field.

Subsequent survey work has further systematised the state of the art, particularly with respect to uncertainty modelling, fusion levels, and application driven design choices. A widely cited review by Khaleghi et al. \cite{khaleghi_multisensor_2013} provides a structured taxonomy of multisensor fusion methods and highlights key challenges related to heterogeneity, scalability, and robustness. While these surveys offer essential coverage of traditional fusion paradigms, they primarily focus on probabilistic and signal level techniques and place less emphasis on interpretability and human centred reasoning.

In parallel, advances in machine learning have led to increased interest in data driven and multimodal fusion methods. Surveys in the multimedia and machine learning communities have examined fusion strategies based on feature learning, representation alignment, and end to end neural architectures \cite{atrey_multimodal_2010, baltrusaitis_multimodal_2019}. These approaches have demonstrated strong empirical performance across a range of applications, but often rely on large annotated data sets and may offer limited transparency in how individual sensor contributions and uncertainties are handled.

At the same time, there has been renewed interest in hybrid and soft computing approaches, including fuzzy logic, rule based systems, and neuro-fuzzy models, which explicitly address uncertainty, imprecision, and context dependence in sensor data. Such methods are particularly relevant in scenarios where sensor reliability varies over time, data are sparse or noisy, or interpretability and explainability are essential requirements. Despite their continued practical relevance, fuzzy and hybrid fusion techniques are often dispersed across application specific studies and lack a unified and up to date tutorial treatment that connects theoretical foundations with contemporary use cases.

This tutorial and survey paper is therefore timely and important for several reasons. First, emerging applications in human robot interaction, medical systems, smart manufacturing, and inclusive workplaces increasingly demand fusion frameworks that are not only accurate, but also transparent, adaptable, and capable of reasoning under uncertainty. Second, the growing availability of heterogeneous sensing technologies, wearable devices, and intelligent environments has intensified the need for principled fusion strategies that can integrate multimodal and subjective data sources. Finally, there is a clear need to bridge classical sensor fusion theory with modern trends in adaptive, human centred, and explainable artificial intelligence.

By synthesising foundational concepts, established surveys, and recent methodological developments, this paper aims to provide a coherent tutorial overview of sensor fusion techniques, with particular emphasis on approaches that balance mathematical rigour, practical adaptability, and interpretability. In doing so, it seeks to serve both newcomers to the field and experienced researchers seeking a structured perspective on current challenges and future research directions.

% ============================================
\section{Fundamentals of Sensor Fusion}
% ============================================

\subsection{Definition and Principles}
Sensor fusion refers to the process of integrating data from multiple sensors to produce a more accurate, reliable, and complete set of information.

\subsection{Types of Sensors Used in Sensor Fusion}
Common sensors used in sensor fusion include IMUs, LiDAR, cameras, accelerometers, and more. These sensors provide data in different modalities, such as temporal, spatial, and spectral data.

\subsection{Levels of Sensor Fusion}
Sensor fusion can be categorised into three levels: low-level (raw data fusion), intermediate-level (feature fusion), and high-level (decision fusion).

\begin{figure}
    \centering
    \includegraphics[width=0.5\linewidth]{ChatGPT Image May 24, 2025, 03_36_04 AM.png}
    \caption{Image Fusion: xxx-Susini}
    \label{fig:enter-label}
\end{figure}

\subsection{Challenges in Sensor Fusion}
Common challenges in sensor fusion include noise, sensor calibration, and handling uncertainties in data integration.

% ============================================
\section{Classical Sensor Fusion Techniques}
% ============================================

\subsection{Bayesian Methods \TC{[Themis]}}
Bayesian methods, such as the Kalman filter and its variants, are widely used in sensor fusion. The Kalman filter, in particular, provides an optimal estimate of the system state by minimising the mean square error.

\subsection{Statistical Data Fusion \TC{[Themis]}} 
Statistical methods, such as maximum likelihood estimation (MLE) and weighted averaging, are often employed to combine data from multiple sources.

\subsection{Evidence Theory \GO{[GAB: rewrite]}}

\GO{[GAB: To be adapted, it is a section of a very old paper of mine]}

\GO{In this section, the Evidence Theory framework is described. Such framework can be used for gaining insights on the ongoing situation, thus implementing a Situation Awareness schema able to infer the most likely cause of faults.}

\GO{Evidence Theory \cite{shafer_mathematical_1976,dempster_upperlower_2008,smets_transferable_1990} is a mathematical formalism to handle uncertainty based on some evidences (e.g. sensors providing information). The idea is to evaluate the \emph{belief} and the \emph{plausibility} of a finite set of hypotheses (e.g. the causes that we want to assess).}

\GO{Let a \textit{frame of discernment} $ \Omega = \{ \omega_1, \cdots, \omega_n\}$ be a set of mutually exclusive hypotheses, and suppose that the hypotheses in $ \Omega$ are exhaustive (e.g. we assume that the possible causes of failure for a telecommunication infrastructure are limited).} 

\GO{Consider a set of sensors, each able to raise alarms originated by several causes or hypotheses. The combination of several simultaneous alarms and thus the combination of different hypotheses leads to the definition of measures of belief and plausibility.}

\GO{Let $\gamma_a$ be a subset of $\Omega$; Evidence Theory operates on the \textit{power set} $\Gamma(\Omega)$ that is the set of all possible combinations of hypotheses $\gamma_a$ in $\Omega$:
\begin{equation}
\Gamma(\Omega)  = \{ \gamma_1, \cdots, \gamma_{2^{|\Omega|}}\}.
\end{equation}
Note that the cardinality of the power set is $|\Gamma(\Omega)| = 2^ {|\Omega|}$, since it contains all possible subsets of $\Omega$, including the empty set $\gamma_1 = \emptyset$ and the universal set $\gamma_{2^{|\Omega|}} = \Omega$.}

\GO{Among the other possible approaches, the Transferable Belief Model (TBM) \cite{smets_transferable_1990} proved to be very effective. TBM is based on the definition of \textit{basic belief mass} function: $m : \Gamma(\Omega) \rightarrow [0, 1] $. This function is a map that assigns to each element of the power set a value between 0 and 1. The function, also referred to as \textit{basic belief assignment} (BBA), shall respect the following constraint:
\begin{equation}
\displaystyle\sum_{\gamma_a\subseteq \Gamma(\Omega)} m(\gamma_a) = 1 \quad\mathrm{ with }\quad m(\emptyset) = 0
\label{eq:bba}
\end{equation}
The set of elements $\gamma_a$ in the power set having $m(\gamma_a)\neq 0$ is called the \textit{focal set}.}

\GO{In this framework, the interest is focused on quantifying the confidence of propositions of the form: ``The true value of $\omega_i$ is in $\gamma_a$,'' with $\gamma_a \in \Gamma(\Omega)$. For $\gamma_a \in \Gamma(\Omega)$, $m(\gamma_a)$ is the part of confidence that supports exactly $\gamma_a$. This means that the true value is in the set $\gamma_a$ but, due to lack of further information, does not support any strict subset of $\gamma_a$. 
This is not a probability function, and it does not respect the property of additivity: $m (\gamma_a \cup \gamma_b) \neq m(\gamma_a) +m (\gamma_b).$} 
 
\GO{Each BBA is an atomic element within the TBM. In fact, each sensor, agent or node must be able to assign the BBA values by some subjective assumptions, or through appropriate algorithms that automatically determine the assignment.}

\GO{In the case of different independent information sources, a rule to aggregate the information must be provided.}

\GO{There are many rules of combination in the literature. Among the others, the most widely used are Dempster's rule and Smets' rule.}

\GO{\textit{Dempster's rule of combination} \cite{dempster_upperlower_2008} was the first to be formalised. The Dempster's rule of combination is purely a conjunctive operation. This rule strongly emphasises the agreement between multiple sources and ignores all the conflicting evidence through a normalisation factor, as shown in \eqref{eq:dempster}. This has the effect of attributing null mass to the empty set. So the rule is formalised as:
\begin{equation}
\nonumber \texttt{Dempster} \{ m_i, m_j\} (\emptyset)  = 0 
\end{equation}
while, for all $\gamma_a\in \Gamma(\Omega)$,
\begin{equation}
\texttt{Dempster} \{ m_i, m_j\} (\gamma_a)  = \frac{\displaystyle\sum_{\gamma_b \cap \gamma_c = \gamma_a} m_i(\gamma_b) m_j(\gamma_c)}{1-\displaystyle\sum_{\gamma_b \cap \gamma_c = \emptyset} m_i(\gamma_b) m_j(\gamma_c)}.
\label{eq:dempster}
\end{equation}}

\GO{\textit{Smets' rule of combination} \cite{smets_transferable_1990} is another rule that allows one to express explicitly the contradiction in the TBM, by letting $m(\emptyset)\neq 0$. This combination rule, compared to Dempster's, simply avoids the normalisation while preserving the commutativity and associativity properties. The formalisation is as follows:
\begin{equation}
\texttt{Smets} \{ m_i, m_j\} (\gamma_a)  = m_i(\gamma_a) \otimes m_j(\gamma_a) \qquad\forall\gamma_a\in \Gamma(\Omega)
\end{equation}
where}
\GO{\begin{equation}
m_i(\gamma_a) \otimes m_j(\gamma_a)= \sum_{\gamma_b \cap \gamma_c =\gamma_a} m_i(\gamma_b) m_j(\gamma_c) \qquad\forall\gamma_a\in \Gamma(\Omega).
\label{eq:smets}
\end{equation}}

\GO{The fact that  $m(\emptyset) > 0$ can be explained in two ways: the open world assumption and the quantified conflict. The closed world assumption, made by Dempster, reflects the idea that the frame of discernment must contain the true value. Necessarily, if the closed world assumption is true, then the set of hypotheses must contain all possibilities. Under this interpretation, being $\emptyset$ the complement of $\Omega$, the mass $m(\emptyset)$ represents the case where the truth is not contained in $\Omega$, and, under the closed world assumption, the value of $m(\emptyset)$ must be zero. On the other hand, if the mass of the $\emptyset$ is not zero, then the meaning is that the truth is not included in the considered hypotheses. The second interpretation of $m(\emptyset) > 0$ is that there is some underlying conflict between the sources that are combined in order to produce the BBA. Hence, the mass assigned to $m(\emptyset)$ represents the degree of conflict. In particular, it can be computed as follows:
\begin{equation}
m_i(\emptyset) \otimes m_j(\emptyset) = 1 - \sum_{\gamma_a\in \Gamma(\Omega), \gamma_a\neq \emptyset} \left( m_i(\gamma_a) \otimes m_j(\gamma_a)\right).
\label{eq:smets_empty}
\end{equation}
The main disadvantage of this approach is its computational complexity, due to computing the power set $\Omega$, that is exponential with respect to the number of hypotheses. This limit usually discourages researchers from applying Evidence Theory to real applications. Some methods to reduce the computational complexity are explained in \cite{shafer_mathematical_1976}.}
\GO{This approach has been also proposed in a distributed framework and applied to the Critical Infrastructure Protection field, \cite{foglietta_networked_2012}.}

\GO{Evidence Theory, as discussed in the previous section, heavily relies on the closed world assumption, that is to say it is particularly ineffective if the set of hypotheses fails to adequately represent the reality of interest.}

\GO{To overcome such a limitation there is the need to define adequate ways to refine the model of the reality under exam. In literature, model refinement is usually considered an off-line task, because of the necessary involvement of human decisions.
Conversely, it is very hard in general to provide automatic tuning mechanisms, and this is especially true for online frameworks.}

\GO{In the case of Evidence Theory, when the model is well-defined and the evidences, detected by the sensors, are enough to discriminate situations, the result is a mass distribution that associates the greatest value to an atomic set (e.g. that the different sensors agree on a single cause).}

\GO{Conversely, if the mass is mainly accumulated on a non-atomic set, then the evidences are not enough to discriminate the actual cause or the model has too much uncertainty and hence is not able to function properly. Indeed, if a knowledge model is not able to distinguish between situations, the domain knowledge has to be re-analysed.}

\GO{Another issue is related to the assumption made in the Dempster--Shafer framework: the closed world assumption. According to such a hypothesis, the considered situations have to be exhaustive and mutually exclusive.}

\GO{Smets, with the TBM approach, overcomes this limit including the possibility that the empty set may have a non-zero mass; in this way it is possible to determine to which extent the estimation is contradictory. The event of an empty set with non-zero mass may happen when several combined sources are in conflict, or when the frame of discernment $\Omega$ does not contain all possible situations, thus highlighting modelling errors and that the truth is not in $\Omega$.}

\GO{Hence, one of the following situations may lead to an empty set with non-zero mass:
\begin{itemize}
\item the knowledge model is correct, but there are problems in the evidence-collecting process or the situation has evolved with time;
\item the knowledge model is incorrect, due to situations not modelled or due to a misidentification of the different situations.
\end{itemize}
}

\GO{In the first case, the problem might be the loss of some evidences due to the process of gathering information. The problem might also be an effect of several situations evolving during time. In a given time period $\left[ t_1,t_2\right]$ evidences support a situation $\gamma_i$, then, when the situation evolves during $\left[ t_2, t_3\right]$ the result is $\gamma_j$; as a result the mass of the empty set is increased.}

\GO{In this case, a solution is to transfer the mass of the empty set to the universal set $\Omega$, when the empty set mass has exceeded a pre-defined threshold. In this way, the system is able to change idea and evaluate the changing situation. This method can be applied in real-time contexts and can be easily implemented.}

\GO{In the latter case, the problem cannot be resolved by simply transferring the mass of the empty set to the universal set. In fact, if the model is incorrect, the conflict value will continue to increase. In this case, other approaches are needed. In this paper, we assume that a new hypothesis can be included in the knowledge model.}

\GO{Consider that two sensors have generated alarms simultaneously, producing high contradiction that is not due to a possible change of situation. So a new hypothesis can be introduced in the knowledge model, linked to the onset of the two specific alarms. In the following numeric example, it has not been verified the possibility that the new hypothesis is mutually exclusive with respect to the others, namely that the new hypothesis is not a subset or partly overlaps with those already considered in the knowledge model.}

\subsection{Optimisation Approaches \GO{[GAB: i can do something here]}}
Optimisation techniques like least squares and gradient-based methods are used to refine the fusion process, improving accuracy and performance.

% ============================================
\section{Advanced Sensor Fusion Techniques}
% ============================================

\subsection{Machine Learning Approaches \TC{[Themis]}}
Machine learning offers powerful tools for sensor fusion, particularly in supervised and unsupervised learning paradigms. Popular methods include decision trees, support vector machines, and principal component analysis (PCA).

\subsection{Deep Learning for Sensor Fusion \TC{[Themis]}}
Deep learning techniques, such as convolutional neural networks (CNNs) and recurrent neural networks (RNNs), have been successfully applied to fuse data from multiple sensor modalities, achieving state-of-the-art results in complex tasks.

\subsection{Multi-sensor SLAM \TC{[Themis]}}
Simultaneous localisation and mapping (SLAM) is an essential application of sensor fusion, often relying on probabilistic models and machine learning algorithms to integrate data from various sensors.

\TC{\cite{jao_uwbzupt_2023}}

\subsection{Sensor Fusion based on Multi-criteria Decision Making \GO{[GAB: rewrite, add relevant literature.]}} 

\color{magenta}
[GAB: TBD add state of the art. For now it is just a section on Multi-criteria decision making from a previous paper of mine. I need to add: (i) some context, (ii) the state of the art of using this for sensor fusion; (iii) how to use it.]

In this subsection, we review the Logarithmic Least Square (LLS) method for solving the Analytic Hierarchy Process problem with multiple decision-makers when the information available is incomplete \cite{bozoki_optimality_2019,oliva_opinion_2019}. We consider $g$ decision-makers, each providing a set of pairwise comparisons $\mathcal{A}^{(u)}_{ij}$ that represents an assessment of the relative importance of the $i$-th alternative over the $j$-th one. Clearly, we have that $A_{ji}^{(u)}=1/A_{ij}^{(u)}$ whenever $A_{ij}^{(u)}$ is available.
The pairwise comparisons $A_{ij}^{(u)}$ are usually defined according to {\em Saaty's scale} \cite{saaty_analytic_1987}, which is summarised in Table~\ref{tab:scale}.
\begin{table*}
\small
\centering
\resizebox{0.95\textwidth}{!}{
\begin{tabularx}{\textwidth}{l|X|X}
\hline
\textbf{$A^{(u)}_{ij}$} & \textbf{Definition} & \textbf{Explanation} \\ \hline
1 & Equal importance & Criterion $i$ and criterion $j$ are considered equally important \\
3 & Moderate importance of one over another & Criterion $i$ moderately more important than criterion $j$ \\
5 & Essential or strong importance & Criterion $i$ is strongly favoured with respect to criterion $j$ \\
7 & Very strong importance & Criterion $i$ is strongly favoured with respect to criterion $j$, and its dominance is demonstrated in practice \\
9 & Extreme importance & The evidence favouring criterion $i$ over criterion $j$
 is of the highest possible order of affirmation \\
 2, 4, 6, 8 & Intermediate values between the two adjacent judgements & When compromise is needed
 \end{tabularx}
 }
 \caption{Saaty's scale for AHP.}
 \label{tab:scale}
 \end{table*}
The pairwise comparisons provided by each decision-maker $u$ are arranged in an $r \times r$ ratio matrix $\mathcal{A}^{(u)}$, which corresponds to a possibly disconnected graph $\mathcal{G}^{(u)}=\{\mathcal{V},\mathcal{E}^{(u)}\}$.
In particular, each alternative can be represented by a node $v_i\in \mathcal{V}$, while the availability of a pairwise comparison $A^{(u)}_{ij}$ can be regarded as an edge joining the $i$-th and $j$-th node.
Let us denote by
\[
\widehat{\mathcal{G}}=\left\{\mathcal{V},\bigcup_{u=1}^g \mathcal{E}^{(u)}\right\}
\]
the graph representing the collective information from all decision-makers (i.e. a graph comprising the union of all edges from each decision-maker, with repeated edges allowed). The optimisation problem to solve is:
\begin{equation}
\label{eq:above1}
{\bm x}^*=\underset{{\bm x}\in\mathbb{R}^r_+}{\arg\min}\,\frac{1}{2}\sum_{u=1}^g\sum_{i=1}^r\sum_{j\in\mathcal{N}^{(u)}_i}\left(
\ln\left(\mathcal{A}_{ij}^{(u)}\right)-\ln\left(\frac{x_j}{x_i}\right)\right)^2,
\end{equation}
where $u$ iterates over all experts, $i$ iterates over all alternatives, and for a given expert $u$ and alternative $i$, $j$ iterates over the neighbourhood $\mathcal{N}^{(u)}_i$ of node $i$ according to the graph $\mathcal{G}^{(u)}$, i.e. the alternatives for which a pairwise comparison with the $i$-th one is provided by the $u$-th expert. To simplify solving the above problem, we substitute $y_i=\ln(x_i)$ and solve the equivalent problem:
\begin{equation}
\label{eq:above}
{\bm x}^*=\exp\left(\underset{{\bm y}\in\mathbb{R}^r}{\arg\min}\,\frac{1}{2}\sum_{u=1}^g\sum_{i=1}^r\sum_{j\in\mathcal{N}^{(u)}_i}\left(
\ln\left(\mathcal{A}_{ij}^{(u)}\right)-y_i+y_j
\right)^2\right),
\end{equation}
where $\exp({\bm y})$ is the component-wise exponential. After substitution, the problem becomes an unconstrained convex minimisation problem, ensuring that a global optimal solution ${\bm y}^*$ is found where the gradient of the objective function, evaluated at ${\bm y}^*$, is zero. Specifically, a global optimal solution satisfies:
\begin{equation}
\label{eq:tobesolved0}
\sum_{u=1}^g L^{(u)}{\bm y}^{*}=\sum_{u=1}^gP^{(u)}{\bm 1}_n,
\end{equation}
where $L^{(u)}$ is the Laplacian matrix associated with $\mathcal{G}^{(u)}$ and $P^{(u)}$ is an $r \times r$ matrix containing the logarithms of the nonzero entries of $\mathcal{A}_{ij}^{(u)}$, with $P^{(u)}_{ij}=0$ when $\mathcal{A}_{ij}^{(u)}=0$. Furthermore, $\exp({\bm y}^*)$ is unique up to a scaling factor if and only if $\widehat{\mathcal{G}}$ is connected. More precisely, since the problem is an unconstrained convex minimisation problem, evaluating the derivative of \eqref{eq:above} at zero reveals that the optimal solution ${\bm y}^*$ satisfies \eqref{eq:tobesolved0}.
\color{black}

\subsection{Fuzzy Logic in Sensor Fusion}
\KD{[DELK: write, add relevant literature]}

Fuzzy logic, first introduced by Zadeh in 1965 \cite{zadeh_fuzzy_1965}, extends classical Boolean logic by allowing partial truth values between 0 and 1. In fuzzy set theory, an element $x$ belongs to a fuzzy set $A$ with a certain degree of membership $\mu_A(x) \in [0,1]$, rather than a crisp inclusion or exclusion as in classical sets. Formally, a fuzzy set $A$ on a universe $X$ is defined as
\begin{equation}
A = \{(x, \mu_A(x)) \mid x \in X, \ \mu_A(x) \in [0,1]\}.
\end{equation}
The function $\mu_A(x)$ is called the \emph{membership function}, representing the degree to which $x$ satisfies the property described by $A$. Logical operations are extended accordingly:
\begin{align}
\mu_{A \cap B}(x) &= \min\{\mu_A(x), \mu_B(x)\}, \\
\mu_{A \cup B}(x) &= \max\{\mu_A(x), \mu_B(x)\}, \\
\mu_{\neg A}(x) &= 1 - \mu_A(x).
\end{align}
This mathematical foundation enables reasoning with vague, imprecise, or uncertain information, as encountered in many real-world systems. Over the decades, fuzzy logic has been widely applied in control, decision-making, and data fusion, providing a framework that mimics human-like reasoning under uncertainty \cite{klir_fuzzy_1995}.

Fuzzy logic is a computational approach based on ``degrees of truth'' rather than binary true/false logic. It is particularly useful in handling uncertainty, imprecision, or vagueness in sensor data. Fuzzy logic is used in sensor fusion to address uncertainty, imprecision, and conflicting data from multiple sensors by employing linguistic rules and membership functions to combine and interpret information. 

Let $S_i$ denote the measurement from the $i$-th sensor. Each sensor output is first mapped to a fuzzy set through a membership function $\mu_i(S_i)$ that expresses the degree of confidence or reliability in that measurement:
\begin{equation}
\mu_i(S_i) = f_i(S_i; \theta_i), \quad \mu_i(S_i) \in [0,1],
\end{equation}
where $f_i(\cdot)$ defines the membership function (e.g. triangular, Gaussian) parameterised by $\theta_i$.

A fuzzy rule base can then be used to express expert knowledge, such as:
\[
\text{IF } S_1 \text{ is Reliable AND } S_2 \text{ is Uncertain THEN } W_1 \text{ is High.}
\]
Each rule $R_j$ is evaluated to produce a firing strength $\alpha_j$, typically computed as
\begin{equation}
\alpha_j = \min_k \big(\mu_{A_{jk}}(S_k)\big),
\end{equation}
where $\mu_{A_{jk}}(S_k)$ is the membership value of input $S_k$ in the $j$-th rule’s antecedent.

To obtain adaptive sensor weights $w_i$, singleton (or centroid-reduced consequents in the case of Mamdani inference) consequents are chosen for the weight outputs. Let $b_{ji} \in \mathbb{R}_{\ge 0}$ denote the representative crisp value associated with the consequent of rule $R_j$ for the $i$-th sensor (e.g. the numeric level implied by ``$W_i$ is High''). Then
\begin{equation}
w_i = \frac{\sum_{j=1}^{N_R} \alpha_j \, b_{ji}}{\sum_{j=1}^{N_R} \alpha_j}.
\end{equation}
The final fused estimate is obtained as a weighted combination of sensor readings:
\begin{equation}
\hat{y} = \frac{\sum_{i=1}^{N_S} w_i \, S_i}{\sum_{i=1}^{N_S} w_i}.
\end{equation}
This produces an explicit, rule-driven adaptive weighting scheme, where each $w_i$ is the normalised aggregation of rule activations $\alpha_j$ mapped through their consequent values $b_{ji}$, and $\hat{y}$ is the corresponding weight-normalised fusion of the sensor measurements.

In \cite{fracczak_fuzzylmw_2016}, a data fusion algorithm is presented that employs fuzzy logic techniques for adaptive sensor weighting. The algorithm adjusts weighting factors dynamically based on fuzzy logic conditions, allowing the fusion process to emphasise the most reliable sensor inputs under varying conditions. This adaptive mechanism results in a more stable and accurate measurement signal. The approach was demonstrated in the context of contactless measurement for determining the linear and angular position of an automatic laparoscopic device (Laparoscope sleeve), showing the practical applicability of fuzzy-based fusion in medical robotics and precision measurement systems. Overall, the study demonstrates how fuzzy logic can enhance sensor fusion by introducing context-sensitive adaptability and robustness, particularly in environments characterised by noise and uncertainty.

The design of inclusive workplaces has become increasingly important, particularly in manufacturing environments where high cognitive demands may disadvantage neurodiverse individuals. Conventional task instructions can elevate cognitive load and error rates, reducing efficiency and inclusivity.

In \cite{testa_fuzzyfusion_2025}, a multi-layer fuzzy logic framework is proposed to assess cognitive load during assembly tasks and support workers with dyslexia. The model integrates physiological, environmental, and task-related data, including heart rate, heart rate variability, electrodermal activity, and eye-tracking metrics from wearable sensors. It particularly targets the reading dimension of cognitive load, which is critical for dyslexic individuals who face challenges with text-based instructions.

A controlled study with 18 participants simulated dyslexia conditions under supported and unsupported scenarios. Results showed that the absence of support increased cognitive load and reduced performance, while adaptive aids such as audio prompts effectively reduced cognitive strain.

In summary, the study demonstrates the potential of fuzzy logic for modelling cognitive states in industrial contexts, supporting the development of personalised and inclusive work environments for neurodiverse individuals.

% ============================================
\section{Applications of Sensor Fusion}
% ============================================

\subsection{Robotics and Autonomous Vehicles}
In robotics and autonomous vehicles, sensor fusion enables accurate localisation, navigation, and perception by combining data from LiDAR, cameras, GPS, and other sensors.

\subsection{Healthcare and Wearable Devices} \KD{[DELK: write, add relevant literature]}

Wearable devices have become an integral component of contemporary healthcare, providing continuous physiological and behavioural measurements in daily life. When these heterogeneous data streams are combined through sensor fusion, they offer more reliable and clinically actionable insights than isolated modalities. This is particularly important in free-living conditions, where sensor noise, motion artefacts and environmental variability can compromise single-sensor performance.

Cardiovascular monitoring is one of the most mature examples of this trend. Optical photoplethysmography alone is highly susceptible to artefacts, yet its fusion with accelerometry or electrocardiography can significantly enhance robustness. Sano and Picard demonstrated that multimodal sensing combining heart rate variability, electrodermal activity and motion features enables reliable detection of stress in daily life \cite{sano_stress_2013}. Similarly, fusion-based approaches have been shown to improve heart rate estimation and the detection of cardiovascular anomalies when users are physically active.

Fall detection illustrates another domain where sensor fusion provides clear benefit. Early single-axis accelerometer systems achieved limited specificity. Wu and colleagues showed that combining tri-axial acceleration with gyroscope data significantly improves discrimination between falls and activities of daily living \cite{wu_fall_2015}. Noury et al.\ further investigated multi-sensor fusion in home-based monitoring environments, showing that fusing inertial and contextual signals reduces false alarms and improves reliability in older adults \cite{noury_fall_2008}.

Rehabilitation and movement analysis also benefit from wearable fusion. Picerno's review highlights how integrating inertial units with electromyography or pressure sensors yields richer biomechanical indicators for gait analysis than any single modality \cite{picerno_wearable_2017}. Such multimodal systems enable remote physiotherapy, continuous mobility assessment and objective tracking of recovery progression.

More broadly, wearable fusion is increasingly used for long-term chronic disease management. Gjoreski et al.\ demonstrated that combining accelerometry, heart rate variability and contextual information enhances performance in detecting stress and affective states in free-living environments \cite{gjoreski_stress_2016}. This supports the use of multi-sensor approaches in mental health monitoring and personalised digital health interventions.

Overall, sensor fusion transforms wearables into adaptive measurement systems capable of handling the variability and noise of real-world settings. This enhances their clinical relevance by improving robustness, reducing false alarms and enabling more holistic modelling of physiological and behavioural states.

\subsection{IoT and Smart Cities} \KD{[DELK: write, add relevant literature]}

The convergence of sensor fusion techniques with Internet of Things infrastructure has enabled the emergence of smart cities, where heterogeneous sensor networks work together to monitor and optimise urban environments. Smart cities integrate diverse sensor data streams from distributed devices to enhance municipal services, improve resource allocation, and elevate quality of life for urban populations \cite{zanella_internet_2014, jin_information_2014}. The rapid increase of low cost sensors, advances in wireless communication, and the maturation of data fusion algorithms have enabled cities to transition from isolated monitoring systems to interconnected intelligent ecosystems.

In smart city deployments, sensor fusion synthesises different data from environmental sensors, traffic monitoring systems, surveillance cameras, and infrastructure embedded instrumentation. Intelligent transportation systems leverage fusion of inductive loop detectors, GPS trajectories, camera feeds, and weather data to estimate traffic flow and optimise signal timing \cite{lv_traffic_2015, guerrero-ibanez_adaptive_2018}. Smart energy grids integrate measurements from smart metres, distributed generation sources, and weather forecasts to balance supply and demand whilst minimising emissions \cite{fang_smart_2012}. The heterogeneity of these data sources necessitates sophisticated fusion architectures that can accommodate asynchronous inputs and handle missing data.

Multi level fusion hierarchies have become prevalent in smart city applications, where raw sensor data undergoes progressive refinement through feature level, decision level, and contextual fusion stages \cite{gao_survey_2020}. Edge computing nodes perform preliminary fusion operations to reduce bandwidth and latency, whilst cloud platforms execute intensive algorithms for citywide optimisation. Recent developments in federated learning enable distributed networks to train fusion models whilst preserving data privacy \cite{lim_federated_2020}. The integration of semantic information from knowledge graphs with sensor measurements facilitates context aware fusion, enabling systems to interpret readings within the broader framework of urban dynamics \cite{petrolo_semantic_2017}.

Environmental monitoring demonstrates the value of sensor fusion, where networks of air quality sensors, meteorological stations, and satellite observations are fused to generate high resolution pollution maps \cite{zheng_diagnosing_2013}. Smart water management systems fuse data from pressure sensors, flow metres, and quality monitors to detect leakage and optimise distribution networks \cite{anumalla_review_2019}. The scalability and resilience of these systems remain important as cities expand to encompass millions of connected devices. Challenges include managing computational complexity, ensuring interoperability across platforms, and maintaining functionality despite sensor failures \cite{kimani_cyber_2019}. Machine learning approaches have demonstrated capabilities in learning fusion mappings from data, though their opaque nature presents obstacles to deployment in safety critical applications \cite{guo_deep_2019}. Looking forwards, the integration of emerging technologies promises to enhance sensor fusion in smart cities, with emphasis shifting towards holistic urban operating systems that raise new questions regarding sustainability and societal implications.

\subsection{Industrial Automation} 
In industrial automation, sensor fusion enhances machine vision, predictive maintenance, and quality control by combining data from various sensor sources.

\subsection{Maritime Situation Awareness \GO{[GAB: I can work on this]}} 

% ============================================
\section{State of the Art: Survey of Current Research}
% ============================================

\subsection{Current Trends in Sensor Fusion}
The field of sensor fusion is evolving rapidly, with new methods emerging that integrate artificial intelligence and machine learning to improve fusion processes.
\KD{[DELK: write, add relevant literature]}

\subsection{Benchmark Datasets and Evaluation Methods}
Common datasets used in sensor fusion research include KITTI, nuScenes, and other benchmarks. Evaluation metrics such as accuracy, precision, recall, and F1-score are used to measure performance.
\KD{[DELK: write, add relevant literature]}

\subsection{Open Challenges and Future Directions}
Despite significant advancements, challenges such as real-time processing, scalability, and the integration of heterogeneous data sources remain. Future research will focus on overcoming these hurdles.
\GO{[GAB]}
\TC{[Themis]}

% ============================================
\section{Tutorial: A Step-by-Step Guide to Implementing Sensor Fusion}
% ============================================

\KD{[DELK: write, add relevant literature]}

This section provides a hands-on tutorial for implementing sensor fusion algorithms, progressing from basic software setup to increasingly advanced fusion strategies. The emphasis is on practical deployment on edge computing platforms and FPGA based systems, which impose strict constraints on latency, computational resources, and determinism. Each example highlights both algorithmic considerations and implementation aspects relevant to real world embedded systems.

\subsection{Basic Setup}

This subsection outlines the basic software and hardware setup required for implementing sensor fusion on edge and embedded platforms. At the software level, Python is widely used for rapid prototyping and algorithm development, while the Robot Operating System (ROS and ROS2) provides a modular middleware for sensor interfacing, message passing, and system integration. Machine learning frameworks such as TensorFlow and PyTorch are commonly employed for learning based fusion approaches, particularly on GPU enabled edge devices.

On platforms such as NVIDIA Jetson Nano, Xavier, and Orin, the typical deployment architecture separates perception and learning tasks from estimation and control. GPU resources are leveraged for image processing and deep neural networks, while CPU cores handle state estimation, filtering, and decision logic to ensure predictable timing behaviour. Hardware in the loop simulation is often used during development to evaluate timing, robustness, and fault handling prior to real world deployment.

Beyond algorithmic implementation, practical sensor fusion systems require integrated development and validation environments that combine sensing, perception, and decision making under realistic operating conditions. In the context of autonomous systems and smart city applications, such environments must support heterogeneous sensors, complex traffic interactions, and human factors.

An example of such a system level validation approach is presented in \cite{deliparaschos_preliminary_2020}, where an integrated testing and simulation infrastructure is developed to support multisensor autonomous vehicle research. Although the work does not focus on a specific fusion algorithm, it highlights the importance of end to end validation frameworks in which sensor fusion components are evaluated alongside perception, control, and human interaction models. This type of infrastructure is particularly relevant for edge and embedded deployments, where sensor fusion performance must be assessed under realistic operational constraints.

For FPGA based systems, toolchains such as Vivado and the Vitis unified development environment, which integrates high level synthesis (HLS) capabilities, enable hardware software co design by allowing computationally intensive estimation algorithms to be mapped to programmable logic. FPGA implementations are particularly attractive when strict real time constraints and low latency are required.

\subsection{Example 1: Implementing a Kalman Filter for Fusion of IMU and GPS Data}

A classical and widely used example of sensor fusion is the integration of inertial measurement unit (IMU) and Global Positioning System (GPS) data using a Kalman filter. IMU sensors provide high rate motion information but suffer from drift, while GPS offers absolute position updates at lower rates and with intermittent availability.

In this example, a discrete time Kalman filter is implemented to estimate the position and velocity of a mobile platform. The prediction step uses IMU acceleration measurements, while GPS position updates are incorporated during the correction step when available. Such architectures are well suited to real time execution on edge devices, as the computational structure is fixed and efficient.

Practical deployments must account for intermittent measurements and non ideal noise characteristics. Recent work on onboard quadrotor navigation demonstrates how Kalman filter based fusion can be extended to handle intermittent GPS like measurements and asynchronous sensor updates in real time \cite{xatziloizou_indoors_2022}. Further robustness can be achieved by adopting alternative estimation criteria, such as the maximum correntropy criterion, which improves performance under heavy tailed noise and outliers \cite{xatziloizou_maximum_2023}.

For FPGA based implementations, Kalman filtering forms a core building block of many fusion pipelines. Efficient hardware software co design approaches enable real time execution while accounting for quantisation effects and fixed point arithmetic. An FPGA based implementation of Kalman and LQG estimation using a hardware in the loop framework illustrates how such estimation cores can be validated and accelerated on programmable logic \cite{deliparaschos_fpgabased_2016}.

\subsection{Example 2: Feature-Level Fusion Using Deep Learning}

Feature level fusion combines intermediate representations extracted from multiple sensors before performing inference or decision making. A common example is the fusion of camera and LiDAR data for perception tasks in autonomous systems. Convolutional neural networks are typically used to extract features from images, while point cloud encoders process LiDAR data.

On edge platforms, feature level fusion is often implemented using GPU acceleration, with neural networks optimised using TensorRT or similar inference frameworks. Architectures such as visual inertial odometry systems demonstrate that tightly coupled fusion of vision and inertial data can be achieved in real time using efficient optimisation based formulations \cite{mourikis_multistate_2007, leutenegger_keyframe_2015}.

Although deep learning based fusion offers strong performance, it introduces challenges related to explainability, data requirements, and computational load. These considerations are particularly important on edge devices, where power and thermal constraints may limit model complexity.

\subsection{Example 3: Decision-Level Fusion in Medical Applications}

Decision level fusion operates on the outputs of independent processing pipelines, combining decisions rather than raw data or features. This approach is particularly suitable in medical and healthcare applications, where heterogeneous sensors measure different physiological phenomena and interpretability is critical.

A representative example is heart rate monitoring using both electrocardiogram (ECG) signals and accelerometer data. ECG provides direct cardiac information, while accelerometers capture motion artefacts and activity levels. Independent classifiers or estimators process each signal stream, and their outputs are combined using logical rules, probabilistic reasoning, or confidence weighted schemes to produce a final decision.

Decision level fusion reduces coupling between sensing modalities and allows individual components to be validated independently, which is often a regulatory requirement in medical systems. Such architectures are well suited to edge and embedded deployment, as they enable modularity and fault isolation.

\textcolor{blue}{[GAB: an additional example could be included here based on expert driven fusion methods, such as the Analytic Hierarchy Process, or evidence theoretic approaches such as Dempster Shafer theory, to illustrate decision level fusion under uncertainty and limited data.]}



% ============================================
\section{Conclusion}
% ============================================

\subsection{Summary of Key Points}
This paper reviewed the fundamentals, advanced techniques, and applications of sensor fusion, providing both a tutorial and survey of the state-of-the-art methods.

\subsection{Implications for Research and Industry}
Sensor fusion will continue to play a key role in various industries, from autonomous systems to healthcare, as technologies evolve and new methods are developed.

\subsection{Future Outlook}
Future advancements in artificial intelligence and computational methods are expected to drive the next wave of innovation in sensor fusion.

\bibliographystyle{IEEEtran}
\bibliography{references}

\end{document}
